\documentclass[10pt,a4paper]{amsart}
\usepackage[margin=19mm]{geometry}
\usepackage{amsmath,amssymb,mathtools}
\usepackage[scr=rsfs,cal=euler]{mathalfa}
\usepackage{xfrac}
\usepackage[unicode,hidelinks,bookmarks=false]{hyperref}
\newcommand{\ie}{i.e.\ }
\newcommand{\cf}{cf.\ }
\newcommand{\resp}{respectively\ }
\newcommand{\Subsection}{\subsection}
\providecommand{\nobreakdash}{\nobreak\hbox{-}}
\setlength{\emergencystretch}{2em}
\multlinegap0pt
\usepackage{amssymb}
\usepackage{tikz}
\newtheorem{theorem}{Theorem}
\newtheorem{remark}{Remark}
\newtheorem{corollary}{Corollary}
\title{Farey Structure in Modulo Krinkle Tilings:\\ Mediant Splicing and Generation of Prototiles from a Single Edge}
\author{Mikihiro Fujiwara}
\date{Source: arXiv:2609.01270v1 (CC BY 4.0)}
\begin{document}
\maketitle

\noindent\emph{Source and licence.} Farey Structure in Modulo Krinkle Tilings:\\ Mediant Splicing and Generation of Prototiles from a Single Edge. Version arXiv:2609.01270v1 (CC BY 4.0), distributed by arXiv under the Creative Commons Attribution 4.0 International licence, http://creativecommons.org/licenses/by/4.0/, as recorded in the arXiv licence metadata for this exact version and shown on the abstract page https://arxiv.org/abs/2609.01270v1. Full licence notice: \texttt{licenses/arxiv-2609.01270-CC-BY-4.0.txt}.\par\smallskip

\noindent\textit{Editorial scope.} Counted unit: the article's theorem thm:decompose (source0.tex lines 1132--1152). The article's complete original proof of it is delivered with it (source0.tex lines 1154--1204, one \texttt{proof} environment of 51 lines). No mathematics is written, repaired or re-derived: the statement and the proof are the article's own lines, in the article's own notation, and no retained line is substituted or re-flowed.  Retained uncounted as the source-local context the proof needs: lines 955--957; lines 1001--1014; lines 1206--1213; lines 1215--1223.  Nothing was omitted from any retained span, so the package carries no omission record.  No retained line contains a citation, so the wrapper carries no bibliography and no bibliography entry.  No retained line contains a figure environment, an image-inclusion command or drawing code, so no image is delivered and the package declares no asset dependency.  Editorial formatting only, in addition: an AMS \texttt{amsart} preamble that restates the article's own declarations and macros used by the retained lines, namely the environments \texttt{theorem}, \texttt{remark}, \texttt{corollary} on separate counters, together with the standard packages those lines need.  The retained environments print in this document as Theorem 1, Remark 1, Corollary 1 in order of appearance, whereas the article prints the same items with the numbering of its own full document.  The complete native source of the article as distributed in the arXiv e-print for this exact version is bundled unchanged as \texttt{sources/209092/source0.tex}.
\begin{equation}
k_1 m - m_1 k \;=\; m_2k_1-m_1k_2 \;=\; 1. \tag{$\star$}\label{eq:star}
\end{equation}

\begin{theorem}[Mediant splicing]\label{thm:splice}
Let $m_1/k_1<m_2/k_2$ be Farey-adjacent, $t\ge2$ common, and $n=t(k_1+k_2)$.
Then the lower path of the $(m_1+m_2,\,k_1+k_2)$-prototile coincides, vertex by vertex,
with the path obtained by concatenating from the origin:
\begin{itemize}
\item the edges $j=0,\dots,k_1-1$ of the lower path of $(m_1,k_1)$, edge $j$ rotated by
$+\,j\cdot 2\pi/(k_1 n)$;
\item the edges $i=0,\dots,k_2-1$ of the lower path of $(m_2,k_2)$, edge $i$ rotated by
$+\,2\pi/n - i\cdot 2\pi/(k_2 n)$;
\item the terminal edge (angle $2\pi/t$, identical in all three tilings).
\end{itemize}
All edges keep unit length; only their directions change, by an angle affine in the
edge number (a \emph{fan-twist}).
\end{theorem}

\begin{remark}[Degenerate blocks]\label{rem:seeds}
A block with $k_i'=1$ is a single unit edge of angle $0$ under either seed, $(0,1)$ or
$(1,1)$, so the geometry cannot distinguish the two seeds and
Theorem~\ref{thm:decompose}(2) cannot force $m_i'$. Requiring the pair to be
Farey-adjacent---equivalently $m_1'+m_2'=m$---selects the canonical value
$m_i'=m_i\in\{0,1\}$; this is the convention used in Corollary~\ref{cor:generation}
below.
\end{remark}

\begin{corollary}[Generation from a single edge]\label{cor:generation}
Iterating Theorem~\ref{thm:decompose} descends the Stern--Brocot tree and terminates,
after finitely many steps, at the seeds $(0,1)$ and $(1,1)$, whose lower words are both
$\ell=(0,1)$: geometrically one and the same unit rhombus (degenerate for $t=2$), and,
at the level of cores, a single unit edge of angle $0$. Reading the recursion upward,
every prototile of the family is obtained from one unit edge by iterated fan-twist
splicing. The leaves of the decomposition tree, read from left to right with
$(0,1)\mapsto 0$ and $(1,1)\mapsto 1$, spell the carry word $c(m,k)$.
\end{corollary}

\begin{theorem}[Unique decomposition]\label{thm:decompose}
Let $(m,k)$ be reduced with $k\ge2$ and $n=tk$, and define the \emph{canonical pair}
\[
k_1 = m^{-1} \bmod k, \qquad m_1 = \frac{mk_1-1}{k}, \qquad
(m_2,k_2)=(m-m_1,\ k-k_1).
\]
\begin{enumerate}
\item[\emph{(1)}] \emph{(Farey parents.)} The canonical pair is Farey-adjacent, it is
the unique Farey-adjacent pair with mediant $m/k$, and its fan-twist splice
(Theorem~\ref{thm:splice}) reproduces the prototile of $(m,k)$.
\item[\emph{(2)}] \emph{(Rigidity.)} Conversely, suppose that for some cut position
$1\le k_1'\le k-1$ and some reduced parameters $(m_1',k_1')$, $(m_2',k_2')$ with
$k_2'=k-k_1'$ (seeds allowed), the fan-twist splice of Theorem~\ref{thm:splice}---with
$(k_1,k_2)$ replaced by $(k_1',k_2')$---reproduces the lower path of $(m,k)$. Then
$k_1'=k_1$, and $m_i'=m_i$ whenever $k_i'\ge2$ (a block with $k_i'=1$ is realized by
either seed; see Remark~\ref{rem:seeds}).
\end{enumerate}
In particular, the prototile admits exactly one fan-twist splice decomposition, and a
fan-twist splice of two prototiles of the family reproduces a prototile of the family
\emph{only if} the parameters are Farey-adjacent.
\end{theorem}

\begin{proof}
(1) $m_1$ is an integer by the definition of the modular inverse, and
\[
m_2k_1-m_1k_2=(m-m_1)k_1-m_1(k-k_1)=k_1m-m_1k=1,
\]
so the pair is Farey-adjacent; $\gcd(m_i,k_i)=1$ follows from the determinant being $1$,
and Theorem~\ref{thm:splice} applied to the pair reproduces the prototile. For
uniqueness among Farey-adjacent pairs: if $(m_1,k_1),(m_2,k_2)$ are Farey-adjacent with
mediant $m/k$, then \eqref{eq:star} gives $k_1m\equiv 1 \pmod k$; since $0<k_1<k$,
necessarily $k_1=m^{-1}\bmod k$, and then $m_1=(mk_1-1)/k$ is forced.

(2) All edges are unit, so the splice reproduces the lower path if and only if it matches
every edge direction. Writing $s'^{(1)},s'^{(2)}$ for the direction indices of
$(m_1',k_1'),(m_2',k_2')$, the angle-matching conditions read
\[
\frac{2\pi s_j}{n}=\frac{2\pi s'^{(1)}_j}{t k_1'}+j\cdot\frac{2\pi}{k_1' n}
\quad(0\le j< k_1'),
\qquad
\frac{2\pi s_{k_1'+i}}{n}=\frac{2\pi s'^{(2)}_i}{t k_2'}+\frac{2\pi}{n}
-i\cdot\frac{2\pi}{k_2' n}
\quad(0\le i< k_2'),
\]
i.e., clearing denominators with $n=tk$,
\begin{equation}
k\,s'^{(1)}_j = k_1'\,s_j - j,
\qquad
k\,s'^{(2)}_i = k_2'\,(s_{k_1'+i}-1) + i.
\tag{$\dagger$}\label{eq:rigid}
\end{equation}
In particular $k$ must divide both right-hand sides.

\emph{The cut position is forced.} Take the last edge of the second block,
$i=k_2'-1$, so that $k_1'+i=k-1$ and $s_{k-1}=k-m$. The right-hand side of the second
equation in \eqref{eq:rigid} is
\[
k_2'(k-m-1)+(k_2'-1) \;=\; k_2'(k-m)-1 \;\equiv\; -(k_2'\,m+1) \pmod k,
\]
so divisibility forces $k_2'\,m\equiv-1\pmod k$. With $1\le k_2'\le k-1$ this gives
$k_2'=k-(m^{-1}\bmod k)$, i.e.\ $k_1'=m^{-1}\bmod k=k_1$.

\emph{The parameters are forced.} If $k_1\ge2$, take $j=1$ in \eqref{eq:rigid}: since
$s_1=m$, we get $s'^{(1)}_1=(k_1m-1)/k=m_1$, and $s'^{(1)}_1=m_1'\bmod k_1'=m_1'$,
whence $m_1'=m_1$. If $k_2\ge2$ then $m\le k-2$ (otherwise $m=k-1$ and $k_2=1$), so
$s_{k_1+1}=m+1$, and $i=1$ in \eqref{eq:rigid} gives, using \eqref{eq:star},
\[
k\,s'^{(2)}_1 = k_2\,m+1 = km-(k_1m-1) = k\,(m-m_1),
\]
whence $m_2'=s'^{(2)}_1=m-m_1=m_2$. Finally, a block with $k_i'=1$ has direction
sequence $(0)$ for both seeds, so \eqref{eq:rigid} imposes no further constraint on that
block (Remark~\ref{rem:seeds}).
\end{proof}

\end{document}
